\section{拒绝预测与容量约束把概率变成行动}
\label{sec:11-Mathematical-Statistics/07-Statistical-Decision-and-Reliable-Prediction/06}

阈值按第二节的公式定成 \(0.0025\)，上线第一天进来 \(263\) 条待复核工单。审核团队六个人，每人每天四十条，能力上限二百。第二天积压 \(63\) 条，第三天 \(131\) 条，两周后队列里最老的工单已经过期，反欺诈的时效性一并作废。

阈值本身没算错，它只是回答了``单看这一条该不该拦''。现实里还有第二个问题：这一批里只能处理二百条，选哪二百条。这一节把处置能力写成约束，接到前面的损失最小化框架上，并把``不作答''正式列为一个行动。

\subsection{拒绝是第三个行动，它的区间由代价定}

设二分类的真实概率 \(p=P(Y=1\given X=x)\)，判错记损失 \(1\)，判对记 \(0\)，转人工记固定代价 \(c\)（\(0<c<1/2\)）。三个行动的条件风险是

\[
\rho(1\given x)=1-p,\qquad
\rho(0\given x)=p,\qquad
\rho(\text{转人工}\given x)=c .
\]

逐点取最小，得到

\[
\begin{cases}
\text{判 }0, & p\le c,\\
\text{转人工}, & c<p<1-c,\\
\text{判 }1, & p\ge 1-c .
\end{cases}
\]

拒绝带的两条边界来自复核成本，不来自``概率接近 \(1/2\) 就该谨慎''这种手感。复核越贵，\(c\) 越大，带子越窄；两类错误代价不对称时，两条边界各自平移，带子不再关于 \(1/2\) 对称。\(c\ge 1/2\) 时带子消失，此时人工复核比直接猜还贵，该做的是提升复核效率而不是接入模型。

\subsection{选择性风险要连覆盖率一起报}

记 \(g(x)\in\set{0,1}\) 表示这一条是否交给模型自动处置。自动覆盖率与被接管样本上的选择性风险分别是

\[
\phi=P\bigl(g(X)=1\bigr),\qquad
R_{\mathrm{sel}}=\E\!\left[L\bigl(Y,\widehat Y\bigr)\given g(X)=1\right].
\]

这两个数必须成对出现。把覆盖率压到 \(30\%\)，只留下最容易的三成，选择性风险自然漂亮，剩下七成全压给人工，系统等于没上。反过来只报``自动化率 \(95\%\)''，被拒绝的那 \(5\%\) 是否集中在某类用户身上就看不见了。可交付的形式是风险--覆盖曲线：横轴覆盖率从 \(1\) 扫到 \(0\)，纵轴是对应的选择性风险，运营方按人力预算在曲线上挑工作点。

拒绝分数用什么也需要论证。Softmax 的最大概率在分布外输入上照样可以很高，模型没见过的东西不会自动表现为``不确定''；集成模型的分歧只反映所选模型族与训练扰动这两个来源。相对稳的做法是用上一节的预测集合大小当拒绝分数——集合里留下三个以上类别就转人工，这个分数背后有覆盖保证撑着，而且集合变大与数据变陌生是同一个方向。无论用哪个分数，拒绝机制本身要在独立数据上单独验证，并检查漂移之后被拒绝的样本是否换了一批人。

\subsection{容量一进来，逐点阈值就变成批内排序}

急诊留观区同时有 \(n\) 位患者，加强监护床位只有 \(B\) 张。设 \(p_i\) 是第 \(i\) 位在普通观察下发生严重恶化的概率，\(a_i\in\set{0,1}\) 表示是否给床位。先假设监护不改变恶化本身，只改变能否及时发现：漏掉一位会恶化的患者损失 \(c_{\mathrm m}\)，占用一张事后证明不需要的床位损失 \(c_{\mathrm o}\)。这一批的期望损失与约束是

\[
\min_{a\in\set{0,1}^n}\ \sum_{i=1}^{n}\left[c_{\mathrm m}p_i(1-a_i)+c_{\mathrm o}(1-p_i)a_i\right],
\qquad \sum_{i=1}^{n}a_i\le B .
\]

把第 \(i\) 位从普通观察改成加强监护，期望损失下降

\[
\Delta_i=c_{\mathrm m}p_i-c_{\mathrm o}(1-p_i).
\]

各人的贡献互不影响，约束只限制个数，于是最优解是把 \(\Delta_i\) 从大到小排，取 \(\Delta_i>0\) 的前至多 \(B\) 位。

\begin{center}
\begin{tikzpicture}[font=\footnotesize,>=stealth,x=0.44cm,y=1.55cm]
  \foreach \i/\h in {1/0.95,2/0.86,3/0.78,4/0.70,5/0.62,6/0.55,7/0.48,8/0.42,
                     9/0.35,10/0.29,11/0.23,12/0.17,13/0.11,14/0.05}
    \fill[black!55] (\i-0.32,0) rectangle (\i+0.32,\h);
  \foreach \i/\h in {1/0.95,2/0.86,3/0.78,4/0.70,5/0.62,6/0.55,7/0.48,8/0.42}
    \fill[black!85] (\i-0.32,0) rectangle (\i+0.32,\h);
  \foreach \i/\h in {15/-0.02,16/-0.09,17/-0.17,18/-0.26,19/-0.36,20/-0.48}
    \fill[black!25] (\i-0.32,0) rectangle (\i+0.32,\h);
  \draw[->,black!70] (0,0) -- (21.4,0);
  \node[black!70,anchor=west] at (16.4,0.42) {名次（\(\Delta_i\) 降序）};
  \draw[black!80,thick] (8.6,-0.62) -- (8.6,1.12);
  \node[black!80,anchor=south] at (8.6,1.12) {容量 \(B=8\)};
  \draw[black!45,dashed] (14.5,-0.62) -- (14.5,0.72);
  \node[black!60,anchor=south west] at (13.2,0.72) {\(\Delta_i=0\)};
  \node[black!60,anchor=north] at (11.5,-0.66) {净收益为正却排不上队};
  \node[black!85,anchor=east,rotate=90] at (-0.35,0.5) {\(\Delta_i\)};
\end{tikzpicture}

\emph{每条工单的净收益排成一列，两条线各管一段：虚线右边送审是亏的，实线右边送审是没人手。当天执行的规则由更靠左的那条线定。}
\end{center}

图里第 \(9\) 到第 \(14\) 位是这套系统的痛处：他们的净收益为正，处置了有赚，但队伍排到第八位就没人了。这段被截掉的收益正是容量约束的影子价格——多加一张床、多雇一个审核员，能挽回的期望损失就是第 \(9\) 位的 \(\Delta_9\)。这个数可以直接拿去谈预算，比``模型准确率提升了两个点''好用得多，它的一般形式见\hyperref[sec:09-Convex-Optimization/05-Duality/04]{``对偶变量是约束的影子价格''}。

排序的对象是净收益而不是概率。两者一致只在一种情形下成立：所有对象的两类代价都相同。现实里复核一条大额转账要半小时、一条小额只要三分钟，挽回的金额也差着量级，此时按 \(p_i\) 排序会把审核员的时间填给一堆廉价工单。更贴近实际的做法是按单位成本收益 \(\Delta_i/\tau_i\) 排序，\(\tau_i\) 是这一条的预计处理时长，约束改成总工时 \(\sum_i \tau_i a_i\le T\)。这是背包问题的连续松弛，贪心按性价比取件在松弛意义下最优，离散最优解与它至多差一件。

还有一个容易被忽略的后果：容量固定时，有效阈值随批次浮动。欺诈高发的那天，第 \(200\) 名的分数会很高，实际执行的阈值远高于 \(0.0025\)；平静的那天则相反。系统实际运行在``取前 \(B\) 名''而不是``过某条线''的规则上，两者只有在流量与基准率稳定时才近似等价。监控要同时记录当天的截断分位点，否则会把基准率的波动误读成模型漂移。

\subsection{排在最前的未必是最该处理的}

前面整套算法把 \(p_i\) 当作``不干预时会出事的概率''。若这个行动本身会改变结果，该排序的量就换了一个：

\[
\E\bigl[Y(1)-Y(0)\given X=x\bigr],
\]

即处理带来的结果改变量，而不是未处理时的风险水平。风险最高的那批人可能无论是否加强监护都会恶化，可改变的空间接近零；中等风险的一批反而能被一次及时处置扭转。按风险排序的资源分配在这种结构下会把床位发给最不可能受益的人，而账面上看起来完全合理——高风险组的恶化率确实最高。这一层需要\hyperref[ch:120]{``因果推断：预测世界不等于改变世界''}提供的工具，本单元只负责说明：给定状态的不确定性和损失怎样选行动。缺了那一层，所谓最优策略可能只是把已有的资源分配重新描述了一遍。

顺带说，top-\(B\) 规则最大化的是总期望收益，它不附带公平性、最低服务率或最坏个体风险的任何承诺。这些要求若真的存在，得写成约束（例如每个分支机构保底 \(10\%\) 的复核名额）或写成多目标，指望概率排序自己满足是没有依据的。第三节的最差子群风险在这里就是可以直接接上的形式。

\subsection{今天的行动会改变明天的数据}

以上都在优化一次行动。被复核的工单会拿到标签，没被复核的不会——下一版模型的训练数据由今天的排序决定，未送审的那部分永远缺标签，模型会越来越确信自己已经学会的那部分，边界处的样本再也不出现。缓解手段是留出小比例的随机抽样名额专门去补标签，用一点即时收益换取长期的可辨识性。这个交换的正式形式在 bandit 与在线学习里叫探索，代价用 regret 计量；若当前行动还改变未来状态本身，单步的后验风险不再等于长期价值，就要交给\hyperref[ch:135]{``最优控制：从反馈稳定到全局最优''}那套以状态、转移与累计回报为骨架的语言。

把这一单元收成一条可执行的检查链：先确认概率对应的是部署人群而非训练人群；再用可靠性图确认这个数字的字面含义可信；接着把两类错误的代价比问出来，算出阈值；需要覆盖承诺就用留出的校准集造集合，并连宽度一起验收；然后把处置能力写成约束，按单位成本净收益排序取前 \(B\) 名，记录当天的截断分位点；最后追问行动会不会改变结果，以及今天的排序会不会污染明天的数据。模型输出的那个概率位于这条链的中间，它前面有人群和校准，后面有代价、容量与反馈。

下一章回到把变异拆开的手艺，从经典线性回归开始，看系数估计怎样接上抽样分布。
